\documentclass{article}
\usepackage[utf8]{inputenc}
\usepackage[T1]{fontenc}
\usepackage{amsmath}
\usepackage{amssymb}
\usepackage{graphicx}
\usepackage{booktabs}
\usepackage{tikz-qtree}
\usepackage{forest}
\usepackage{gb4e}
\usepackage{expex}
\usepackage{hyperref}
\usepackage{url}

\graphicspath{{./figures/}}

\title{Syntactic Structure and Constituent Ordering in Natural Language}

\author{Author Name}

\begin{document}

\maketitle

\begin{abstract}
This paper investigates syntactic structures in natural language using formal
grammars and treebank analyses. We examine constituent order variations
across languages and their implications for linguistic theory.
\endabstract}

\section{Introduction}

The study of syntactic structure is fundamental to understanding how natural
languages encode meaning through grammatical constructions.

\section{Syntax Trees}

Syntactic trees provide a hierarchical representation of sentence structure:

\ex Standard sentence structure
\gll John loves Mary .
John loves Mary .
\glft `John loves Mary.'
\endgl
\xe

The corresponding tree structure:

\begin{exe}
\ex
\begin{forest}
[S
  [NP [Det [The]] [N [cat]]]
  [VP [V [chased]] [NP [Det [the]] [N [dog]]]]]
\end{forest}
\xe

Using \textbf{tikz-qtree}:

\begin{center}
\Tree [.S [.NP [.Det The ] [.N cat ] ]
           [.VP [.V chased ]
               [.NP [.Det the ] [.N dog ] ] ] ]
\end{center}

\section{Glossed Examples}

Interlinear glossed text (IGT) is standard in linguistic publications:

\ex Multiple morpheme example
\gll wola-ka-$\emptyset$-a
want-PRES-DECL-IND
\glft `S/he wants it.'
\xe

\ex Complex verbal complex
\gll chi-whelu-yu
  1SUBJ-see-PAST
\glft `I saw (him/her/it).'
\xe

\section{Binding Theory}

\ex Binding Condition A
\gll John$_i$ criticized himself$_i$.
John criticize-PAST REFL
\glft `John criticized himself.'
\xe

\ex Binding Condition B
\gll He$_i$ criticized John$_i$.
3SG criticize-PAST John
\glft `He criticized John.'
\xe

\section{Conclusion}

This paper has examined syntactic structures using modern linguistic methodology.

\bibliographystyle{plain}
\bibliography{refs}

\end{document}
